\documentclass[10pt,a4paper]{amsart}
\usepackage[margin=19mm]{geometry}
\usepackage{amsmath,amssymb,mathtools}
\usepackage[scr=rsfs,cal=euler]{mathalfa}
\usepackage{xfrac}
\usepackage[unicode,hidelinks,bookmarks=false]{hyperref}
\newcommand{\ie}{i.e.\ }
\newcommand{\cf}{cf.\ }
\newcommand{\resp}{respectively\ }
\newcommand{\Subsection}{\subsection}
\providecommand{\nobreakdash}{\nobreak\hbox{-}}
\setlength{\emergencystretch}{2em}
\multlinegap0pt
\usepackage{amssymb}
\usepackage{dsfont}
\usepackage{esint}
\usepackage[all]{xy}
\usepackage{tikz}
\newtheorem{mlem}{Lemma}
\newtheorem{lemma}{Lemma}
\newcommand{\R}{\mathbb{R}}
\newcommand{\mb}{\mbox}
\title{Interior regularity for fully fractional evolution equations and a priori estimates on unbounded domains}
\author{Wenxiong Chen, Yahong Guo, Congming Li}
\date{Source: arXiv:2502.07530v2 (CC BY 4.0)}
\begin{document}
\maketitle

\noindent\emph{Source and licence.} Interior regularity for fully fractional evolution equations and a priori estimates on unbounded domains. Version arXiv:2502.07530v2 (CC BY 4.0), distributed by arXiv under the Creative Commons Attribution 4.0 International licence, http://creativecommons.org/licenses/by/4.0/, as recorded in the arXiv licence metadata for this exact version and shown on the abstract page https://arxiv.org/abs/2502.07530v2. Full licence notice: \texttt{licenses/arxiv-2502.07530-CC-BY-4.0.txt}.\par\smallskip

\noindent\textit{Editorial scope.} Counted unit: the article's lemma thm2.1.2 (source0.tex lines 850--855). The article's complete original proof of it is delivered with it (source0.tex lines 858--993, one \texttt{proof} environment of 136 lines). No mathematics is written, repaired or re-derived: the statement and the proof are the article's own lines, in the article's own notation, and no retained line is substituted or re-flowed.  Retained uncounted as the source-local context the proof needs: lines 399--408; lines 783--797; lines 813--815.  Nothing was omitted from any retained span, so the package carries no omission record.  No retained line contains a citation, so the wrapper carries no bibliography and no bibliography entry.  No retained line contains a figure environment, an image-inclusion command or drawing code, so no image is delivered and the package declares no asset dependency.  Editorial formatting only, in addition: an AMS \texttt{amsart} preamble that restates the article's own declarations and macros used by the retained lines, namely the environments \texttt{mlem}, \texttt{lemma} on separate counters, together with the standard packages those lines need.  The retained environments print in this document as Mlem 1, Lemma 1 in order of appearance, whereas the article prints the same items with the numbering of its own full document.  The complete native source of the article as distributed in the arXiv e-print for this exact version is bundled unchanged as \texttt{sources/173112/source0.tex}.
The above approach can be interpreted as a {\em directional perturbation average}. At a given point $x$, it is impossible to control the derivative of $v$ solely by the value of $v$ itself. However, after making directional perturbations from $x$ to $x^j = x + \eta_j$ for $ j=1,2, \cdots, 2^n$ along directions $\eta_j$, and then taking the average, we are able to realize such a control. Such an estimate is essentially performed on the fully fractional heat kernel $G(x-y, t-\tau)$. For convenience of future applications, we summarized this as
follows:
\begin{mlem}[ Fully  fractional heat kernel estimates ]\label{key0} In each quadrant $I_j$ of $R^n$, there is a vector $\eta_j\in\partial B_{\frac{1}{\sqrt{n}}}(0)$ along the diagonal direction such that
\begin{equation}\label{key1} G(y,\tau)\leq {\bf{e^{-\frac{|y|}{n\tau}}}}\sum\limits_{j=1}^{2^n}G(y+{\eta}_j,\tau),  \forall \tau>0, y\in B_1^c(0).\end{equation}
Furthermore,   for each $\tau>0, y\in B_1^c(0)$, it holds
\begin{equation} D_y^\alpha  G(y,\tau)\leq C\sum\limits_{j=1}^{2^n}G(y+{\eta}_j,\tau), \forall \, |\alpha|=1 \, or \, 2,\end{equation}
and
 \begin{equation}  \partial_\tau G(y,\tau)\leq C\sum\limits_{j=1}^{2^n}G(y+{\eta}_j,\tau),\end{equation}
here $C=C(n)$ is a universal constant.
\end{mlem}

Subsequently, for any fixed $t\in (1,2)$ and $\tau<t,$ we obtain
\begin{eqnarray}\label{est key}
\frac{|y-x|}{t-\tau} e^{-\frac{|y-x|^2}{4(t-\tau)}}
&\leq & \frac{|y-x|}{t-\tau} e^{-\frac{|y-x^j|^2+\delta^2|y-x|}{4(t-\tau)}} \nonumber \\
& \leq & \frac{|y-x|}{t-\tau} e^{-\delta^2\frac{|y-x|}{(t-\tau)}}e^{-\frac{|y-x^j|^2}{4(t-\tau)}}\nonumber\\
& \leq & Ce^{-\frac{|y-x^j|^2}{4(t-\tau)}},
\end{eqnarray}
here, we have used the basic inequality
$$X^p\leq  C_\delta e^{\delta^2X}, \ \forall X\geq 0, p>0.$$
 Furthermore, we can also derive a general type inequality as follows,
\begin{eqnarray}\label{key est2}
\left(\frac{|y-x|}{t-\tau}\right)^p e^{-\frac{|y-x|^2}{4(t-\tau)}}
\leq & Ce^{-\frac{|y-x^j|^2}{4(t-\tau)}}, \  \forall y\in I_j\cap B^c_2(0), \tau<t, \forall p>0.
\end{eqnarray}
The above arguments also validates Lemma \ref{key0}.

\begin{equation}\label{est-vij}
|v_{x_{i}x_{j}}(x,t)|\leq C_{n,s}\int_{-\infty}^t \int_{\mathbb{R}^n}  \frac{f_{Q^c}(y,\tau)}{(t-\tau)^{n/2+1-s}} \left[\frac{1}{2(t-\tau)} +\frac{|x-y|^2}{4(t-\tau)^2}\right]e^{-\frac{|x-y|^2}{4(t-\tau)}}dy d\tau:=I_1+I_2.
\end{equation}

\begin{lemma} \label{thm2.1.2}
There exists a constant $C$ independent of $x\in B_1(0),$ such that
$$
\|v(x,\cdot)\|_{C^{\alpha}_{t} ((1,2))} \leq C \|v(x,\cdot)\|_{L^\infty ((0,3))},\ \ \forall \alpha\in(0,1).
$$
\end{lemma}

\begin{proof} Fixed $x\in B_1(0),$ for any $\varphi\in C_0^\infty((1,2)),$
\begin{eqnarray}\label{est-t}
\begin{aligned}
\int_1^2\frac{\partial v(x,t)}{\partial t}\varphi(t)dt&=-\int_1^2v(x,t)\varphi'(t)dt\\
&=-\int_1^2v(x,t)\left(\lim\limits_{\delta\to 0}\frac{\varphi(t)-\varphi(t-\delta)}{\delta}\right)dt\\
&=-\lim\limits_{\delta\to 0}\frac{1}{\delta}\left\{\int_1^2v(x,t)\varphi(t)dt-\int_1^2v(x,t)\varphi(t-\delta)dt\right\}\\
&=-\lim\limits_{\delta\to 0}\frac{1}{\delta}\left\{\int_1^2v(x,t)\varphi(t)dt-\int_{1-\delta}^{2-\delta}v(x,t+\delta)\varphi(t)dt\right\}\\
&=-\lim\limits_{\delta\to 0}\frac{1}{\delta}\left\{\int_1^2[v(x,t)-v(x,t+\delta)]\varphi(t)dt+\left(\int_{2-\delta}^{2}
-\int_{1-\delta}^{1}\right)v(x,t+\delta)\varphi(t)dt\right\}\\
&=-\lim\limits_{\delta\to 0}\int_1^2\frac{v(x,t)-v(x,t+\delta)}{\delta}\varphi(t)dt.
\end{aligned}
\end{eqnarray}
Here we have utilized the fact that
\begin{equation}\label{small-term}\lim\limits_{\delta\to 0}\frac{1}{\delta}\int_{2-\delta}^{2}v(x,t+\delta)\varphi(t)dt=\lim\limits_{\delta\to 0}\frac{1}{\delta}\int_{1-\delta}^{1}v(x,t+\delta)\varphi(t)dt=0.\end{equation}

Indeed, by the {\em mean value theorem for integrals},   the boundedness of $v(x,t)$ and the fact that $\varphi\in C_0^\infty((1,2)),$ we have
\[\lim\limits_{\delta\to 0}\frac{1}{\delta}\int_{2-\delta}^{2}v(x,t+\delta)\varphi(t)dt=\lim\limits_{\xi\to 2}v(x,\xi)\varphi(\xi)=0.\]
Similarly,
\[\lim\limits_{\delta\to 0}\frac{1}{\delta}\int_{1-\delta}^{1}v(x,t+\delta)\varphi(t)dt=\lim\limits_{\xi\to 1}v(x,\xi)\varphi(\xi)=0.\]
These imply \eqref{small-term}.

Now for any fixed $(x,t)\in B_1(0)\times(1,2):=\tilde{Q}$, substituting the expression
\begin{equation}v(x,t)=\int_{-\infty}^t \int_{\mathbb{R}^n} (f_{Q^c}(y,\tau) G(x-y, t-\tau) dy d\tau\end{equation}
into the right hand side  of \eqref{est-t}, we derive
\begin{eqnarray}\label{est-t1}
\begin{aligned}
&\int_1^2\frac{v(x,t+\delta)-v(x,t)}{\delta}\varphi(t)dt\\
=&\frac{1}{\delta}\int_1^2\left\{\int_{-\infty}^{t+\delta}\int_{\R^n}f_{Q^c}(y,\tau)G(x-y,t+\delta-\tau)dyd\tau
-\int_{-\infty}^{t}\int_{\R^n}f_{Q^c}(y,\tau)G(x-y,t-\tau)dyd\tau\right\}\varphi(t)dt\\
=&\int_1^2\int_{-\infty}^{t}\int_{\R^n}f_{Q^c}(y,\tau)\frac{G(x-y,t+\delta-\tau)-G(x-y,t-\tau)}{\delta}dyd\tau\varphi(t)dt\\
&+\frac{1}{\delta}\int_1^2\int_{t}^{t+\delta}\int_{\R^n}f_{Q^c}(y,\tau)G(x-y,t+\delta-\tau)dyd\tau\varphi(t)dt\\
:=&J_1+J_2.\end{aligned}
\end{eqnarray}

First, we want show that $J_1$  is equivalent to the differentiation of $v(x,t)$ with respect to $t$ within the integral sign, that is,
\begin{eqnarray}\label{J1-lim}
\begin{aligned}
J_1\to&\int_1^2\int_{-\infty}^{t}\int_{\R^n}f_{Q^c}(y,\tau)\frac{\partial G}{\partial t} (x-y,t-\tau)dyd\tau\varphi(t)dt, \ \mb{as}\ \ \delta\to 0.
\end{aligned}
\end{eqnarray}

 In fact, applying the Cauchy mean value theorem, we obtain
\begin{eqnarray*}
\begin{aligned}J_1=\int_1^2\int_{-\infty}^{t}\int_{\R^n}f_{Q^c}(y,\tau)\frac{\partial G}{\partial t}(x-y,t+\delta'-\tau)dyd\tau\varphi(t)dt, \delta'\in(0,\delta).\end{aligned}
\end{eqnarray*}
Differentiate $G(x,t)$ with respect to $t$,
\begin{equation}\label{grn-t}\frac{\partial G(x,t)}{\partial t}%=\frac{\partial}{\partial_\tau}\left(\frac{C_{n,s}}{{\tau}^{n/2+1-s}} e^{-\frac{|y|^2}{4\tau}}\right)
=\frac{C_{n,s}}{{t}^{n/2+1-s}} e^{-\frac{|x|^2}{4t}}(\frac{|x|^2}{4t^2}-\frac{\frac{n}{2}+1-s}{t}).\end{equation}
Then for $\tau<t, 0<\delta'<\delta,$
\begin{eqnarray}\label{delt}
\begin{aligned}\left|\frac{\partial G}{\partial t}(x-y,t+\delta'-\tau)\right|\leq \frac{C e^{-\frac{|x-y|^2}{4(t+\delta'-\tau)}}}{{(t+\delta'-\tau)}^{n/2+1-s}}\left(\frac{|x-y|^2}{(t+\delta'-\tau)^2}+\frac{1}{t+\delta'-\tau}\right)\end{aligned}
\end{eqnarray}

Similar to the argument as in deriving the estimate of $v_{x_i}$, we estimate the right hand side of  \eqref{delt} in two possible cases.

Case 1. $y\in B_2^c(0),$ then $|y-x|\geq 1.$ We can choose $x^j$ as above such that for each $j=1,2,...2^n,$ if $y\in I_j,$ by \eqref{key est2}(with $p=1,2$), it holds
\begin{eqnarray}\label{delt1}
\begin{aligned}e^{-\frac{|x-y|^2}{4(t+\delta'-\tau)}}\left(\frac{|x-y|^2}{(t+\delta'-\tau)^2}+\frac{1}{t+\delta'-\tau}\right)\leq Ce^{-\frac{|x^j-y|^2}{4(t+\delta'-\tau)}}.
\end{aligned}
\end{eqnarray}

Case 2. $y\in B_2(0), \tau\leq 0,$ then
\[|x-y|\leq 3, t+\delta'-\tau\geq1.\]
Therefore,
\begin{eqnarray}\label{delt2}
\begin{aligned}e^{-\frac{|x-y|^2}{4(t+\delta'-\tau)}}\left(\frac{|x-y|^2}{(t+\delta'-\tau)^2}+\frac{1}{t+\delta'-\tau}\right)\leq Ce^{-\frac{|x-y|^2}{4(t+\delta'-\tau)}}.
\end{aligned}
\end{eqnarray}

Substituting \eqref{delt1} and \eqref{delt2} into \eqref{delt}, we deduce
\begin{eqnarray}\label{delt3}
\begin{aligned}\left|\frac{\partial G}{\partial t}(x-y,t+\delta'-\tau)\right|&\leq C\sum\limits_{j=0}^{2^n}\frac{e^{-\frac{|x^j-y|^2}{4(t+\delta'-\tau)}}}{{(t+\delta'-\tau)}^{n/2+1-s}}\\
&\leq  C\sum\limits_{j=0}^{2^n} G(x^j-y,t+\delta'-\tau)\\
&\leq  C\sum\limits_{j=0}^{2^n}\{ G(x^j-y,t-\tau)+ G(x^j-y,t+1-\tau)\},\end{aligned}
\end{eqnarray}
%��t-\tau\to0��\infty����case�ֱ���Ƶġ�
{with} $x^0=x$.
In addition,
\begin{eqnarray}\label{delt4}
\begin{aligned}&\int_1^2\int_{-\infty}^{t}\int_{\R^n}f_{Q^c}(y,\tau)\sum\limits_{j=0}^{2^n}\{ G(x^j-y,t-\tau)+ G(x^j-y,t+1-\tau)\}dyd\tau\varphi(t)dt\\
&\leq C\left(\sum\limits_{j=0}^{2^n}\int_1^2{v(x^j,t+1)\varphi(t)dt+\int_1^2v(x^j,t)}\varphi(t)dt\right)\\
&\leq C \|v\|_{L^{\infty}(Q)}\leq C.\end{aligned}
\end{eqnarray}

Therefore, taking into account of \eqref{delt3}, \eqref{delt4}, and by the {\em Lebesgue's dominated convergence theorem}, we verify \eqref{J1-lim}.  This implies that, as $\delta\to 0,$
\begin{eqnarray}\label{est-J1}
\begin{aligned}
J_1\to&\int_1^2\int_{-\infty}^{t}\int_{\R^n}\frac{C_{n,s}f_{Q^c}(y,\tau)}{{(t-\tau)}^{n/2+1-s}} e^{-\frac{|x-y|^2}{4(t-\tau)}}\left[\frac{|x-y|^2}{4(t-\tau)^2}-\frac{\frac{n}{2}+1-s}{t-\tau}\right]dyd\tau\varphi(t)dt.\\
\end{aligned}
\end{eqnarray}

Next we claim \begin{equation}\label{est-J2}J_2\to 0, \ \mb{as}\ \ \delta\to 0.\end{equation}

Indeed, for any $0<\delta<1,$ if $\tau\in(0,\delta)$, then $\tau+t\in(0,3)$, by $supp f_{Q^c}\subset Q^c$,
\[D:=\{y\in\R^n:f_{Q^c}(y,t+\tau)\neq 0\}\subset B_2^c(0).\]
 Now for $\forall y\in D$, $|y-x|\geq 1$,  by equality \eqref{grn-t}, we observe that $G(|x-y|,\delta)$ is monotone increasing in $\delta$ for $\delta\ll1$.  % because
%\[\frac{\partial G(y,\tau)}{\partial_\tau}%=\frac{\partial}{\partial_\tau}\left(\frac{C_{n,s}}{{\tau}^{n/2+1-s}} e^{-\frac{|y|^2}{4\tau}}\right)
%=\frac{C_{n,s}}{{\tau}^{n/2+2-s}} e^{-\frac{|y|^2}{4\tau}}(\frac{|y|^2}{16\tau}-(\frac{n}{2}+1-s))\geq 0.\]
Thus, for sufficiently small $\delta>0,$ we obtain
\begin{eqnarray}\label{est-t2}
\begin{aligned}
J_2&=\frac{1}{\delta}\int_1^2\int_{t}^{t+\delta}\int_{\R^n}f_{Q^c}(y,\tau)G(x-y,t+\delta-\tau)dyd\tau\varphi(t)dt\\
&=\frac{1}{\delta}\int_1^2\int_{0}^{\delta}\int_{\R^n}f_{Q^c}(y,\tau+t)G(x-y,\delta-\tau)dyd\tau\varphi(t)dt\\
&\leq\frac{1}{\delta}\int_1^2\int_{0}^{\delta}\int_{\R^n}f_{Q^c}(y,\tau+t)G(x-y,\delta)dyd\tau\varphi(t)dt\\
&=\frac{1}{\delta}\int_{0}^{\delta}\int_{\R^n}\int_1^2f_{Q^c}(y,\tau+t)G(x-y,\delta)\varphi(t)dtdyd\tau\\
&=\frac{1}{\delta}\int_{0}^{\delta}\int_{\R^n}\int_{1+\tau}^{2+\tau}f_{Q^c}(y,t)G(x-y,\delta)\varphi(t-\tau)dtdyd\tau\\
&\leq \sup_{t\in(1,2)}\varphi(t)\frac{1}{\delta}\int_{0}^{\delta}\int_{\R^n}\int_{1}^{\frac{5}{2}}f_{Q^c}(y,t)G(x-y,\delta)dtdyd\tau\\
&\leq C\int_{\R^n}\int_{1}^{\frac{5}{2}}f_{Q^c}(y,t)G(x-y,\delta)dtdy:=J_3.\\
\end{aligned}
\end{eqnarray}
Therefore, it suffices to prove
\begin{equation}\label{est-J3}J_3=C\int_{\R^n}\int_{1}^{\frac{5}{2}}f_{Q^c}(y,\tau)G(x-y,\delta)d\tau dy\to 0, \ \mb{as}\ \ \delta\to 0.\end{equation}

In fact,
\[f_{Q^c}(y,\tau)G(x-y,\delta)=f_{Q^c}(y,\tau)\frac{e^{-\frac{|x-y|^2}{4\delta}}}{{\delta}^{n/2+1-s}}\to 0 \ a.e.\ \mb{as}\ \ \delta\to 0. \]
 Furthermore, there exists a point $t_0\in (\frac{5}{2},3)$ such that for any $\tau\in(1,\frac{5}{2})$ and sufficiently small $\delta>0$, it holds
 \[f_{Q^c}(y,\tau)G(x-y,\delta)\leq f_{Q^c}(y,\tau)G(x-y,t_0-\tau)\]
 and \begin{eqnarray*}
\begin{aligned}
&\int_{\R^n}\int_1^3f_{Q^c}(y,\tau)G(x-y,t_0-\tau)d\tau dy\\
\leq &\int_{\R^n}\int_{-\infty}^{t_0}f_{Q^c}(y,\tau)G(x-y,t_0-\tau)d\tau dy\\
=&v(x,t_0)\leq \|v\|_{L^{\infty}(Q)}.\end{aligned}
\end{eqnarray*}
Hence, by the {\em Lebesgue's dominated convergence theorem}, \eqref{est-J3} is valid.

Together \eqref{est-t2} with \eqref{est-J3}, we verified claim \eqref{est-J2}.

Then combining \eqref{est-t}, \eqref{est-t1},\eqref{est-J1} and \eqref{est-J2}, we arrive at
\[\int_1^2\frac{\partial v(x,t)}{\partial t}\varphi(t)dt=C_{n,s}\int_1^2\int_{-\infty}^{t}\int_{\R^n}\frac{f_{Q^c}(y,\tau)}{{(t-\tau)}^{n/2+1-s}} e^{-\frac{|x-y|^2}{4(t-\tau)}}\left[\frac{|x-y|^2}{4(t-\tau)^2}-\frac{\frac{n}{2}+1-s}{t-\tau}\right]dyd\tau\varphi(t)dt,\]
which implies that
\[\frac{\partial v(x,t)}{\partial t}=C_{n,s}\int_{-\infty}^{t}\int_{\R^n}\frac{f_{Q^c}(y,\tau)}{{(t-\tau)}^{n/2+1-s}} e^{-\frac{|x-y|^2}{4(t-\tau)}}\left[\frac{|x-y|^2}{4(t-\tau)^2}-\frac{\frac{n}{2}+1-s}{t-\tau}\right]dyd\tau, a.e. (x,t)\in \tilde{Q}.\]
Now similar to the argument for the estimate of $v_{x_ix_j}$ in \eqref{est-vij}, we deduce that the weak derivative of $v(x,t)$ with respect to the variable $t$ is controlled by $\|v\|_{L^{\infty}(Q)}$, that is,
\[\Bigg|\frac{\partial v(x,t)}{\partial t}\Bigg|\leq C\|v\|_{L^{\infty}(Q)}.\]
Therefore, for any fixed $x\in B_1(0),$ $v(x,\cdot)\in W^{1,\infty}((1,2)),$  by the sobolev  embedding theorem, we deduce that $v(x,\cdot)\in C_t^{0,1}((1,2))\subset C_t^{\alpha}((1,2)),  \, \forall \alpha\in(0, 1),$ and
\[\|v(x,\cdot)\|_{C_t^{\alpha}((1,2))}\leq C \|v(x,\cdot)\|_{L^{\infty}((0,3))}.\]
This completes the proof of Lemma \ref{thm2.1.2}. \end{proof}

\end{document}
